% !TEX program = xelatex
\documentclass[11pt,a4paper]{article}

% ---------------------------
% Packages
% ---------------------------
\usepackage[margin=1in]{geometry}
\usepackage{kotex}
\usepackage{amsmath, amssymb, bm}
\usepackage{graphicx}
\usepackage{booktabs}
\usepackage{longtable}
\usepackage{array}
\usepackage{float}
\usepackage{caption}
\usepackage{subcaption}
\usepackage{natbib}
\usepackage{hyperref}
\usepackage{xcolor}
\usepackage{setspace}
\usepackage{enumitem}
\usepackage{url}

\hypersetup{
    colorlinks=true,
    linkcolor=blue!50!black,
    citecolor=blue!50!black,
    urlcolor=blue!50!black
}

\captionsetup{font=small, labelfont=bf}
\setlength{\parindent}{1.2em}
\setlength{\parskip}{0.35em}
\onehalfspacing

% If a figure file is not uploaded yet, show a placeholder instead of failing.
\newcommand{\maybeincludegraphics}[2][]{%
  \IfFileExists{#2}{\includegraphics[#1]{#2}}{%
    \fbox{\parbox[c][0.22\textheight][c]{0.88\linewidth}{\centering
    Figure file not found: \\[0.5em]
    \texttt{\detokenize{#2}} \\[0.5em]
    Run \texttt{make\_overleaf\_figures.R} and upload the generated \texttt{figures/} folder to Overleaf.}}
  }%
}

\title{Path of Exile 일중 동시접속자 곡선의 구조 분해와 함수형 이상치 탐지}
\author{통계학과 신현수 \quad 2025712931}
\date{2026년 6월}

\begin{document}
\maketitle

\begin{abstract}
본 연구는 Steam 공개 동시접속자 데이터를 이용하여 \textit{Path of Exile}의 일중(intraday) 동접 곡선을 함수형 데이터로 분석하였다. 각 날짜의 24시간 동접 시계열을 하나의 함수로 보고, B-spline 평활화, 함수형 주성분 분석(FPCA), function-on-scalar regression, 함수형 이상치 탐지를 적용하였다. 분석 결과, 일중 곡선의 변동은 거의 단일 모드로 설명되었으며, 첫 번째 함수형 주성분이 전체 변동의 94.47\%를 설명하였다. 이 모드는 곡선의 모양을 유지한 채 전체 수준만 위아래로 이동시키는 수준 효과(level effect)로 해석되었다. 두 번째 함수형 주성분은 리그 출시일에만 나타나는 단기적 이벤트 효과를 포착하였다. 또한 주말, 리그 출시일, 리그 출시 후 첫 주 및 첫 달의 효과를 시간대별 함수로 추정한 결과, 리그 출시일은 출시 전 점검 시간의 음의 효과와 출시 후 저녁 시간대의 폭발적 양의 효과를 동시에 갖는 매우 비대칭적인 패턴을 보였다. Outliergram과 MS plot을 결합한 함수형 이상치 탐지에서는 강력한 이상치 93개를 확인하였고, 분석 가능한 모든 리그 출시일이 이상치로 검출되었다. 본 연구는 함수형 데이터 분석이 게임 서비스의 수요 구조, 이벤트 효과, 비정상 날짜 패턴을 분해하고 해석하는 데 유용함을 보여준다.
\end{abstract}

\noindent\textbf{주요어:} 함수형 데이터 분석, 온라인 게임, 동시접속자 데이터, 함수형 주성분 분석, 함수형 회귀, 함수형 이상치 탐지

\newpage
\tableofcontents
\newpage

% =========================================================
\section{서론}
% =========================================================

\subsection{연구 배경}
함수형 데이터 분석(Functional Data Analysis; FDA)은 시간이나 공간과 같이 연속적인 인덱스를 따라 관측되는 곡선 자료를 하나의 통계적 객체로 다루는 방법론이다 \citep{ramsay2005functional}. 게임 서비스의 동시접속자 수는 함수형 분석의 좋은 대상이다. 하루 24시간이라는 자연스러운 정의역을 가지며, 고밀도 시계열로 관측되고, 패치, 이벤트, 점검, 주말 등의 요인이 곡선의 수준뿐 아니라 모양에도 영향을 미치기 때문이다.

일반적인 시계열 분석은 특정 시점의 값이나 일별 총량에 초점을 두는 경우가 많다. 그러나 게임 서비스의 실제 운영에서는 단순히 동접이 증가했는지보다 \emph{언제} 증가하고 \emph{어느 시간대}가 민감한지가 중요하다. 예를 들어 리그 출시일은 단순히 동접이 많은 날이 아니라, 출시 전 점검과 출시 후 폭증이라는 비대칭적인 일중 구조를 가질 수 있다. 이러한 시간대별 구조를 포착하기 위해서는 날짜별 일중 시계열을 함수로 다루는 접근이 자연스럽다.

\subsection{연구 대상}
본 연구는 \textit{Path of Exile}(이하 PoE)의 Steam 공개 동시접속자 데이터를 사용한다. PoE는 약 3개월 주기의 리그(league) 출시로 잘 알려져 있으며, 새 리그 출시는 동접의 급격한 증가를, 리그 후반으로 갈수록 점진적인 감소를 동반한다. 이처럼 뚜렷한 이벤트 구조는 FPCA, function-on-scalar regression, 이상치 탐지와 같은 함수형 분석 기법을 적용하기에 적합한 환경을 제공한다.

\subsection{연구 질문}
본 연구는 다음 세 가지 질문을 다룬다.
\begin{enumerate}[label=(\arabic*)]
    \item PoE의 일중 동접 곡선 변동은 소수의 주요 모드로 어떻게 분해되는가?
    \item 주말, 리그 출시일, 리그 출시 후 첫 주 및 첫 달의 효과는 어느 시간대에서 크게 나타나는가?
    \item 비정상적인 일별 곡선은 어떻게 식별되고, 어떤 유형으로 해석될 수 있는가?
\end{enumerate}

% =========================================================
\section{데이터}
% =========================================================

\subsection{데이터 출처와 구성}
본 연구는 Mendeley Data에 공개된 \textit{Steam Games Dataset: Player count history, Price history and data about games}를 사용하였다 \citep{wannigamage2020steam}. 이 데이터셋은 Steam 플랫폼에서 운영되는 게임들의 동시접속자 수를 5분 간격으로 기록한 시계열을 제공한다. 본 연구에서는 PoE(Steam appid: 238960)에 해당하는 데이터를 추출하여 분석하였다.

\begin{table}[H]
\centering
\caption{데이터 개요}
\begin{tabular}{ll}
\toprule
항목 & 내용 \\
\midrule
게임 & Path of Exile \\
Steam appid & 238960 \\
분석 기간 & 2017-12-14 -- 2020-08-12 \\
최종 일수 & 971일 \\
관측 해상도 & 5분 간격 \\
일일 관측점 수 & 288개 \\
원시 관측치 수 & 280,224개 \\
\bottomrule
\end{tabular}
\end{table}

원시 데이터는 UTC 기준으로 저장되어 있었다. PoE는 특정 지역에 국한되지 않은 글로벌 게임이므로, 분석 전반에서 UTC를 유지하였다. 이는 글로벌 피크 시간대, 즉 유럽 저녁과 북미 오후가 겹치는 시간대를 자연스럽게 해석하는 데 유리하다.

\subsection{전처리}
전처리는 다음 순서로 수행하였다. 첫째, \texttt{players = 0}인 관측과 결측치를 모두 비정상 관측으로 간주하고 결측 처리하였다. 원자료에서 결측치는 1,790개(0.64\%)였고, 0으로 기록된 관측치는 570개였다. 둘째, 날짜별 결측 개수를 계산하여 50\% 이상 결측을 갖는 2개 날짜(2018-02-10, 2018-02-11)를 분석에서 제외하였다. 셋째, 각 날짜 내부의 결측 관측치는 선형 보간으로 채워 공통 5분 grid 위의 완전한 함수형 행렬을 구성하였다. 마지막으로 분산 안정화를 위해 원시 동접자 수 $y$에 대해 $\log(1+y)$ 변환을 적용하였다.

\subsection{탐색적 시각화}

\begin{figure}[H]
\centering
\maybeincludegraphics[width=0.95\linewidth]{figures/fig01_overall_timeseries.pdf}
\caption{PoE 전체 동시접속자 시계열. 약 3개월 주기의 spike는 리그 출시 효과를 반영한다.}
\label{fig:overall_ts}
\end{figure}

그림 \ref{fig:overall_ts}에서 약 3개월 주기의 spike가 반복적으로 관찰된다. 각 spike는 새 리그 출시일과 대응하며, 리그 시작 직후 동접이 폭증하고 이후 점진적으로 감소하는 톱니파(sawtooth) 구조를 보인다.

\begin{figure}[H]
\centering
\maybeincludegraphics[width=0.95\linewidth]{figures/fig02_spaghetti.pdf}
\caption{날짜별 일중 동접 곡선의 spaghetti plot. 각 곡선은 하루 24시간의 $\log(1+\text{players})$ 변화를 나타낸다.}
\label{fig:spaghetti}
\end{figure}

그림 \ref{fig:spaghetti}의 일별 곡선들은 유사한 일중 패턴을 공유하면서 전체 수준만 상하로 이동하는 형태를 보인다. 이는 변동의 대부분이 모양 변화보다 수준 변화로 설명될 가능성을 시사한다.

\begin{figure}[H]
\centering
\maybeincludegraphics[width=0.98\linewidth]{figures/fig03_mean_variance.pdf}
\caption{평균 일중 곡선과 분산 함수. 평균 곡선은 UTC 06--07시에 최저점, UTC 19--20시에 최고점을 형성한다. 분산은 저녁 피크 시간대에서 증가한다.}
\label{fig:mean_var}
\end{figure}

그림 \ref{fig:mean_var}에서 평균 곡선은 UTC 06--07시에 최저점을 보이고 UTC 19--20시에 최고점을 형성한다. 분산 함수는 저녁 시간대로 갈수록 증가하며, 특히 UTC 17시 이후 변동성이 뚜렷하게 커진다. 이는 평일과 이벤트일의 차이가 주로 글로벌 피크 시간대에서 크게 나타남을 시사한다.

% =========================================================
\section{방법}
% =========================================================

\subsection{B-spline 평활화}
각 날짜의 관측 곡선 $x_i(t)$는 B-spline basis의 선형결합으로 표현하였다.
\begin{equation}
    x_i(t) = \sum_{k=1}^{K} c_{ik}\phi_k(t) = \bm{\phi}(t)^\top \mathbf{c}_i .
\end{equation}
곡선의 거칠기를 제어하기 위해 2차 도함수 roughness penalty를 사용한 penalized least squares를 적용하였다.
\begin{equation}
    \mathrm{PENSSE}_{\lambda}(\mathbf{c}_i)
    = \sum_{j=1}^{288} \{y_{ij} - \bm{\phi}(t_j)^\top \mathbf{c}_i\}^2
    + \lambda \int [D^2 x_i(t)]^2\,dt .
\end{equation}
본 연구에서는 정의역 $[0,24]$에서 51개의 cubic B-spline basis를 사용하였고, GCV를 기준으로 탐색한 결과 $\lambda=10^{-2}$를 평활화 강도로 채택하였다.

\subsection{함수형 주성분 분석}
평활화된 곡선의 변동을 함수형 주성분으로 분해하였다. 평균 함수 $\bar{x}(t)$를 제거한 중심화 곡선을 다음과 같이 표현한다.
\begin{equation}
    x_i(t) - \bar{x}(t) = \sum_{k=1}^{\infty} \xi_{ik}\,\psi_k(t),
\end{equation}
여기서 $\psi_k(t)$는 직교정규한 주성분 함수이고, $\xi_{ik}$는 $i$번째 날짜의 $k$번째 주성분 점수이다. 본 연구에서는 \texttt{fda::pca.fd()}를 사용하여 처음 4개의 함수형 주성분을 추출하였다.

\subsection{Function-on-scalar regression}
주말과 리그 관련 dummy 변수가 날짜별 곡선에 미치는 영향을 보기 위해 function-on-scalar regression을 적용하였다.
\begin{equation}
\begin{aligned}
    y_i(t) =\;& \beta_0(t)
    + \beta_1(t)\,\mathrm{Weekend}_i
    + \beta_2(t)\,\mathrm{LeagueDay}_i \\
    &+ \beta_3(t)\,\mathrm{LeagueWeek1}_i
    + \beta_4(t)\,\mathrm{LeagueMonth1}_i
    + \epsilon_i(t).
\end{aligned}
\end{equation}
각 회귀계수 함수 $\beta_k(t)$ 역시 B-spline basis로 확장하고 roughness penalty를 적용하여 추정하였다. 추정에는 \texttt{fda::fRegress()}를 사용하였다.

\begin{table}[H]
\centering
\caption{Function-on-scalar regression에 사용한 설명 변수}
\begin{tabular}{ll}
\toprule
변수 & 정의 \\
\midrule
Weekend & 토요일 또는 일요일이면 1, 그 외 0 \\
LeagueDay & 리그 출시일이면 1, 그 외 0 \\
LeagueWeek1 & 리그 출시 후 1--6일이면 1, 그 외 0 \\
LeagueMonth1 & 리그 출시 후 7--29일이면 1, 그 외 0 \\
\bottomrule
\end{tabular}
\end{table}

\subsection{함수형 이상치 탐지}
비정상적인 일별 곡선을 식별하기 위해 두 가지 함수형 이상치 탐지 방법을 적용하였다. Outliergram은 modified band depth(MBD)와 modified epigraph index(MEI)를 이용하여 shape outlier를 식별한다 \citep{arribas2014outliergram}. MS plot은 magnitude outlyingness와 variation outlyingness를 동시에 고려하여 수준 이상치와 모양 이상치를 구분한다 \citep{dai2018msplot}. 본 연구에서는 두 방법 중 하나라도 이상치로 탐지한 날짜를 전체 이상치 후보로 보고, 두 방법 모두에서 탐지한 날짜를 강력한 이상치(strong outlier)로 정의하였다.

% =========================================================
\section{결과}
% =========================================================

\subsection{함수형 주성분 분석 결과}

\begin{table}[H]
\centering
\caption{함수형 주성분의 분산 설명 비율}
\begin{tabular}{ccc}
\toprule
FPC & 설명 비율 & 누적 설명 비율 \\
\midrule
FPC1 & 94.47\% & 94.47\% \\
FPC2 & 3.11\% & 97.58\% \\
FPC3 & 0.78\% & 98.36\% \\
FPC4 & 0.37\% & 98.73\% \\
\bottomrule
\end{tabular}
\label{tab:fpca_var}
\end{table}

표 \ref{tab:fpca_var}에 따르면 첫 번째 주성분은 전체 변동의 94.47\%를 설명하고, 두 번째 주성분까지 누적 97.58\%를 설명한다. 이는 PoE 일중 곡선 변동이 사실상 두 개의 주요 모드로 요약될 수 있음을 의미한다.

\begin{figure}[H]
\centering
\maybeincludegraphics[width=0.98\linewidth]{figures/fig04_fpca_pm.pdf}
\caption{평균 곡선에 $\pm 2$ 표준편차의 함수형 주성분을 더한/뺀 형태. FPC1은 전체 수준의 평행 이동, FPC2는 저녁 피크 시간대의 구조 변화를 나타낸다.}
\label{fig:fpca_pm}
\end{figure}

그림 \ref{fig:fpca_pm}에서 FPC1은 평균 곡선을 거의 평행하게 위아래로 이동시키는 수준 효과(level effect)로 해석된다. 반면 FPC2는 저녁 피크의 강도와 위치를 조절하는 모드로, 리그 출시일과 같은 특수 날짜에서 크게 나타난다.

\begin{figure}[H]
\centering
\maybeincludegraphics[width=0.98\linewidth]{figures/fig05_fpca_scores_time.pdf}
\caption{FPC1 및 FPC2 점수의 시계열. FPC1은 리그 사이클 전체를, FPC2는 리그 출시일의 단기 이벤트 효과를 반영한다.}
\label{fig:fpc_scores_time}
\end{figure}

그림 \ref{fig:fpc_scores_time}에서 FPC1 점수는 약 3개월 주기의 톱니파 패턴을 보인다. 반면 FPC2 점수는 대부분의 날짜에서는 0 근처에 머물다가 리그 출시일에만 큰 spike를 보인다. 따라서 FPC1은 ``리그 사이클 효과'', FPC2는 ``출시일 이벤트 효과''로 해석할 수 있다.

\subsection{함수형 회귀 결과}

\begin{figure}[H]
\centering
\maybeincludegraphics[width=0.98\linewidth]{figures/fig06_fosr_effects_ci.pdf}
\caption{주말, 리그 출시일, 리그 첫 주, 리그 첫 달의 시간대별 효과 함수와 95\% 신뢰구간.}
\label{fig:fosr_effects}
\end{figure}

그림 \ref{fig:fosr_effects}에서 주말 효과는 하루 전반에 걸쳐 양수이지만 특히 UTC 09--12시에서 크게 나타난다. 리그 출시일 효과는 가장 극적이다. UTC 17시 이전에는 음수 효과(점검 시간), UTC 19시 이후에는 강한 양수 효과(출시 후 폭증)를 보이며, 하루 전체 평균으로는 설명되지 않는 비대칭 구조를 드러낸다. 리그 첫 주와 첫 달의 효과는 비교적 평행 이동형이며, 첫 주의 효과가 첫 달보다 크다.

\begin{figure}[H]
\centering
\maybeincludegraphics[width=0.82\linewidth]{figures/fig07_timevarying_r2.pdf}
\caption{시간대별 결정계수 $R^2$. 단순 dummy covariate만으로도 일중 곡선 변동의 상당 부분을 설명한다.}
\label{fig:r2}
\end{figure}

시간대별 결정계수의 평균은 60.2\%이며, 최대는 UTC 17.42시 부근, 최소는 UTC 22.50시 부근에서 나타났다. 주말과 리그 메타데이터만으로도 PoE 동접 곡선의 주요 변동을 상당 부분 설명할 수 있음을 보여준다.

\subsection{함수형 이상치 탐지 결과}

\begin{figure}[H]
\centering
\begin{subfigure}{0.48\linewidth}
\centering
\maybeincludegraphics[width=\linewidth]{figures/fig08_outliergram.pdf}
\caption{Outliergram}
\end{subfigure}
\hfill
\begin{subfigure}{0.48\linewidth}
\centering
\maybeincludegraphics[width=\linewidth]{figures/fig09_msplot.pdf}
\caption{MS plot}
\end{subfigure}
\caption{함수형 이상치 탐지 결과. Outliergram은 shape outlier를, MS plot은 magnitude 및 variation outlyingness를 시각화한다.}
\label{fig:outlier_methods}
\end{figure}

Outliergram은 96개, MS plot은 145개의 이상치를 식별하였다. 두 방법에서 동시에 탐지된 강력한 이상치는 93개였다. 특히 분석 가능한 모든 리그 출시일이 강력한 이상치로 탐지되었다.

\begin{table}[H]
\centering
\caption{강력한 이상치의 유형별 분류}
\begin{tabular}{lll}
\toprule
유형 & 개수 & 해석 \\
\midrule
League start day & 10 & 리그 출시일: 점검 후 저녁 폭증 \\
League week 1 & 9 & 리그 출시 후 첫 주의 평행 이동형 증가 \\
Mid-league spike & 9 & 리그 중 이벤트성 spike \\
League fatigue (low) & 2 & 리그 후반 동접 저하 \\
Shape: event-like & 3 & 출시일과 유사한 비출시일 shape \\
Shape: maintenance & 6 & 특정 시간대 점검 dip 패턴 \\
Other & 54 & 기타 약한 이상치 \\
\bottomrule
\end{tabular}
\label{tab:outlier_type}
\end{table}

\begin{figure}[H]
\centering
\maybeincludegraphics[width=0.98\linewidth]{figures/fig10_outlier_types.pdf}
\caption{이상치 유형별 일중 곡선. 회색은 정상 곡선이며, 색상은 이상치 유형을 나타낸다.}
\label{fig:outlier_types}
\end{figure}

그림 \ref{fig:outlier_types}에서 리그 출시일 곡선은 출시 전 점검과 출시 후 폭증이 결합된 비대칭 패턴을, 리그 첫 주 곡선은 정상 곡선의 평행 이동을, 점검일 곡선은 특정 시간대 dip을 뚜렷하게 보인다. 정상 곡선 집합이 매우 좁은 띠를 형성한다는 점도 확인된다.

\begin{figure}[H]
\centering
\maybeincludegraphics[width=0.98\linewidth]{figures/fig11_case_study.pdf}
\caption{대표적 이상치 4개 사례와 평균 곡선의 비교. 점검형, 이벤트형, 피로형 이상치가 서로 다른 구조를 보임을 확인할 수 있다.}
\label{fig:case_study}
\end{figure}

그림 \ref{fig:case_study}는 대표적 이상치 사례를 보여준다. 2020-03-13은 점검 후 저녁 시간대 폭증이 결합된 전형적인 리그 출시일 이상치이며, 2019-10-09는 특정 시간대 sharp dip을 갖는 전형적인 점검형 이상치이다. 반면 2018-12-09는 평균과 거의 같은 모양을 유지한 채 전체 수준만 크게 올라간 평행 이동형 사례다.

% =========================================================
\section{논의}
% =========================================================
본 연구는 PoE의 일중 동시접속자 곡선에 함수형 데이터 분석을 적용하여 세 가지 핵심 결과를 얻었다. 첫째, PoE의 일중 변동은 매우 단순하다. 첫 번째 함수형 주성분 하나가 전체 변동의 94.47\%를 설명하며, 이는 곡선 모양의 변화보다 전체 수준의 변동이 훨씬 크다는 뜻이다. 둘째, 함수형 분석은 시간 척도가 다른 효과를 자연스럽게 분해한다. FPCA에서는 90일에 걸친 리그 사이클 효과와 단 하루의 출시일 효과가 서로 다른 주성분으로 분리되었고, 함수형 회귀에서는 주말 효과, 출시일 효과, 첫 주 및 첫 달 효과가 시간대별 함수로 명시적으로 추정되었다. 셋째, 함수형 이상치 탐지는 게임 운영에서 중요한 날짜를 자동으로 식별한다. 리그 출시일은 모두 이상치로 검출되었고, 점검일과 피로 구간도 다른 유형으로 구분되었다.

본 연구의 한계도 존재한다. Steam 공개 데이터만 사용했기 때문에 점검 일정, 마케팅 활동, 스트리머 방송, 경쟁작 출시 등 외생 요인을 완전히 통제하지 못했다. 또한 분석 대상이 PoE 한 게임이므로 결과의 정량적 규모를 다른 게임에 그대로 일반화하기는 어렵다. 그럼에도 본 연구는 공개 동시접속자 데이터만으로도 게임 서비스의 수요 구조를 함수형 관점에서 분해하고 해석할 수 있음을 보여주었다.

% =========================================================
\section{결론}
% =========================================================
본 연구는 PoE의 일중 동시접속자 곡선을 함수형 데이터로 보고, 평활화, FPCA, 함수형 회귀, 함수형 이상치 탐지를 적용하였다. 그 결과 PoE의 일중 변동은 주로 수준 변화로 설명되며, 리그 사이클과 출시 이벤트가 서로 다른 시간 척도의 변동으로 분리되고, 이상치 탐지는 출시일과 점검일을 효과적으로 식별함을 확인하였다. 이러한 분석 틀은 다른 라이브서비스 게임이나 공개 플랫폼 트래픽 데이터에도 적용 가능하며, 운영 이벤트가 시간대별 이용 구조를 어떻게 바꾸는지 이해하는 데 유용한 도구가 될 수 있다.

\newpage
\bibliographystyle{apalike}
\bibliography{references}

\newpage
\appendix
\section{분석 기간 내 PoE 리그 출시일}
\begin{table}[H]
\centering
\caption{분석 기간 내 PoE 리그 출시일}
\begin{tabular}{lll}
\toprule
순번 & 리그명 & 출시일 \\
\midrule
1 & Abyss & 2017-12-08 (데이터 시작 전) \\
2 & Bestiary & 2018-03-02 \\
3 & Incursion & 2018-06-01 \\
4 & Delve & 2018-08-31 \\
5 & Betrayal & 2018-12-07 \\
6 & Synthesis & 2019-03-08 \\
7 & Legion & 2019-06-07 \\
8 & Blight & 2019-09-06 \\
9 & Metamorph & 2019-12-13 \\
10 & Delirium & 2020-03-13 \\
11 & Harvest & 2020-06-19 \\
\bottomrule
\end{tabular}
\end{table}

\section{분석 환경 요약}
\begin{table}[H]
\centering
\caption{분석 환경 요약}
\begin{tabular}{ll}
\toprule
항목 & 내용 \\
\midrule
분석 대상 기간 & 2017-12-14 -- 2020-08-12 \\
최종 일수 & 971일 \\
일일 관측점 수 & 288개 \\
평활화 basis & 51개 cubic B-spline \\
평활화 $\lambda$ & $10^{-2}$ \\
FPCA 개수 & 4개 주성분 \\
주요 R 패키지 & \texttt{fda}, \texttt{roahd}, \texttt{fdaoutlier}, \texttt{ggplot2} \\
\bottomrule
\end{tabular}
\end{table}

\end{document}
